\chapter{CONCLUSIONES Y PERSPECTIVAS}

Este capítulo sintetiza los aportes de la tesis, responde de forma ordenada a los objetivos planteados en el Capítulo 1 y traza las perspectivas de trabajo futuro. La estructura sigue la lógica de los objetivos, de manera que el lector pueda verificar en qué medida cada uno fue atendido, y concluye con las líneas que quedan abiertas para investigaciones posteriores.

\section{Respuestas a los Objetivos de Investigación}

El primer objetivo específico buscaba evaluar el impacto del desbalance de clases sobre la robustez de las explicaciones estructurales generadas por GNNExplainer, PGExplainer y GNNShap a través de escenarios de desbalance. La evidencia indica que el desbalance no es el factor dominante que la formulación inicial suponía. En el eje Elliptic, la estabilidad presenta un perfil esencialmente plano a lo largo de los escenarios, sin un deterioro localizado, y en el eje sintético el balanceo exhibe un tamaño de efecto despreciable sobre la calidad de las explicaciones. La conclusión es que la degradación de la estabilidad no se explica principalmente por el nivel de desbalance, sino por la elección del explicador, un resultado que reorienta la pregunta original hacia el factor que verdaderamente gobierna la calidad explicativa.

El segundo objetivo buscaba comparar la resiliencia de GCN, GraphSAGE, GAT y TAGCN midiendo la consistencia de sus explicaciones. La respuesta, una vez corregidos dos artefactos de evaluación sucesivos y replicado el entrenamiento con tres semillas de modelo, es que la estabilidad no ordena a las cuatro arquitecturas en cuatro posiciones distinguibles, sino que las separa en dos grupos, uno alto formado por GAT y GCN y uno bajo formado por GraphSAGE y TAGCN, con diferencias significativas entre grupos y no dentro del grupo bajo. La aparente ventaja de GraphSAGE que reportaba una versión anterior no sobrevive a la reparación de la métrica de estabilidad, y la replicación confirma que su lugar está en el grupo bajo. Esta partición se reproduce en los dos regímenes de densidad, ya que el grafo sintético denso separa los mismos dos grupos con las mismas composiciones, de manera que la resiliencia explicativa se manifiesta con relativa independencia de la densidad del grafo. La hipótesis de partida, que anticipaba a TAGCN como la arquitectura más estable por su alcance multisalto, queda refutada, puesto que TAGCN aparece de forma consistente en el grupo bajo en las tres semillas. Conviene precisar el grano de esta respuesta, ya que lo que la evidencia sostiene es la pertenencia a un grupo y no un ordenamiento fino dentro de él, y en el caso de TAGCN su dispersión entre semillas, superior al triple de la de las demás arquitecturas, aconseja no fijar su valor central con excesiva precisión.

El tercer objetivo buscaba analizar el impacto de las estrategias de mitigación del desbalance sobre la interpretación del modelo. La conclusión es que la ponderación de clases y la Focal Loss apenas afectan a la plausibilidad o a la estabilidad de las explicaciones, con tamaños de efecto muy inferiores a los del explicador o la arquitectura. Este resultado tiene una implicación práctica valiosa, ya que sugiere que la elección de la estrategia de balanceo puede realizarse en función del rendimiento predictivo sin temor a distorsionar de forma apreciable la calidad de las explicaciones. La tesis, por tanto, no encuentra evidencia de que balancear los datos fabrique una representación menos fiel del fenómeno ilícito, al menos en lo que respecta a las estrategias basadas en la función de pérdida.

El cuarto objetivo buscaba construir una matriz de recomendación técnica que articulara rendimiento y estabilidad explicativa. La tesis concluye que no existe una tríada de arquitectura, explicador y balanceo que sea uniformemente óptima, porque la estabilidad, la plausibilidad y la fidelidad son dimensiones independientes que ningún método optimiza a la vez. La recomendación resultante es condicional al propósito de la auditoría, de modo que, en cuanto a arquitectura, la elección pertinente es la del grupo alto formado por GAT y GCN, indistintamente, frente al grupo bajo, mientras que entre los explicadores PGExplainer destaca por su plausibilidad cuando el objetivo es reconocer el patrón de lavado, y GNNExplainer por su fidelidad cuando el objetivo es auditar el razonamiento del modelo. Formular la recomendación de arquitectura en términos de grupo y no de una arquitectura concreta es más fiel a la evidencia disponible y, además, más útil en la práctica, porque deja libre la elección dentro del grupo alto para atender a otros criterios como el coste computacional o la disponibilidad de implementaciones. Esta matriz condicional es más útil que una recomendación única, porque refleja la naturaleza multidimensional de la calidad explicativa que el estudio puso de manifiesto.

\section{Aportes Principales}

El aporte central de la tesis es la demostración, respaldada estadísticamente, de que la estabilidad, la plausibilidad y la fidelidad de las explicaciones de GNN son dimensiones independientes en el dominio AML. A este aporte se suman tres hallazgos concretos, la superioridad robusta de PGExplainer en plausibilidad de aristas sobre un grafo con \textit{ground-truth}, la ausencia de un puente entre estabilidad y plausibilidad, y la disociación entre plausibilidad y fidelidad que solo un régimen sintético controlado permite exhibir. En el plano metodológico, la tesis aporta la corrección del subgrafo receptivo, que reordena el ranking de estabilidad y evita confundir un fallo de memoria con una propiedad de la arquitectura, la distinción entre la degeneración de un método y la limitación impuesta por la dispersión del dato, y la distinción entre la reproducibilidad de los pesos de un modelo y la reproducibilidad de sus conclusiones, establecida al reentrenar la matriz completa y comprobar que el cómputo sobre aceleradores gráficos no devuelve modelos idénticos pero sí el mismo filtro de calidad y la misma partición de estabilidad. El diseño de dos ejes complementarios, uno real y otro sintético, constituye en sí mismo un aporte, ya que sostiene conclusiones que ningún dataset por separado permitiría.

Conviene destacar, por último, un aporte de naturaleza transversal, la práctica de la honestidad metodológica que atraviesa el trabajo. La disposición a reconocer las limitaciones del dato, a retractar hallazgos apoyados en artefactos y a distinguir con cuidado entre lo que la evidencia respalda y lo que apenas sugiere, no es un aporte marginal, sino una contribución al modo en que deben conducirse los estudios de explicabilidad. En un campo donde los resultados son sensibles al protocolo y las conclusiones se propagan con rapidez, esta disciplina de contraste y corrección es tan valiosa como los hallazgos concretos que produce, y constituye un ejemplo de cómo la evidencia debe primar sobre las expectativas de partida.

\section{Limitaciones}

Las conclusiones anteriores deben leerse dentro de las fronteras que el estudio reconoce de forma explícita. El eje real se apoya en un único conjunto de datos anonimizado cuyo desplazamiento temporal restringe el análisis de estabilidad a la partición de validación. El análisis inferencial descansa en tres semillas de modelo por celda en ambos ejes, un número suficiente para separar los dos grupos de arquitecturas pero modesto para ordenar dentro de cada grupo, y que resulta particularmente ajustado en el caso de TAGCN sobre Elliptic, cuya dispersión entre semillas es la mayor de las cuatro arquitecturas. El grafo sintético carece del desplazamiento temporal y del ruido de las transacciones reales, de modo que su validez es interna, y la comparación de fidelidad entre familias de explicadores utiliza máscaras distintas por construcción, lo que obliga a interpretarla con cautela. Estas limitaciones no comprometen los hallazgos centrales, pero delimitan su alcance y motivan las perspectivas que se enuncian a continuación.

\section{Perspectivas Futuras}

Varias direcciones de trabajo se desprenden de forma natural de los resultados obtenidos. Una primera línea consiste en ampliar el eje real a conjuntos de datos adicionales de transacciones, preferiblemente con features no anonimizadas, que permitieran medir la plausibilidad también sobre datos reales y contrastar la disociación hallada en el eje sintético. Una segunda línea consiste en enriquecer el grafo sintético con desplazamiento temporal y con tipologías más variadas, de manera que aproxime mejor las condiciones del dato real sin perder el \textit{ground-truth} que lo hace evaluable. Una tercera línea consiste en incrementar el número de semillas de modelo y de realizaciones del grafo para estrechar los intervalos de confianza y afinar los tamaños de efecto, así como en explorar explicadores adicionales más allá de los tres considerados.

Una cuarta línea, de carácter más aplicado, consiste en trasladar la matriz de recomendación condicional a un flujo de auditoría concreto, en el que el propósito de cada revisión determine el explicador empleado, y en evaluar con analistas de cumplimiento si la distinción entre plausibilidad y fidelidad resulta operativa en la práctica. Finalmente, la disociación entre plausibilidad y fidelidad abre una pregunta teórica de fondo, la de si es posible diseñar explicadores que optimicen ambas dimensiones a la vez, o si existe un compromiso fundamental entre alinear una explicación con el conocimiento del dominio y alinearla con el mecanismo del modelo. Responder a esa pregunta excede el alcance de esta tesis, pero constituye una continuación natural de sus hallazgos y una contribución potencial de gran valor para la explicabilidad de las redes neuronales sobre grafos.

\section{Reflexión Final}

Esta tesis partió de una pregunta sobre la estabilidad de las explicaciones bajo desbalance extremo y terminó por reformular el propio marco desde el que se juzga la calidad de una explicación. La evidencia reunida a lo largo de los dos ejes experimentales indica que la calidad explicativa no es una magnitud única, sino un espacio de al menos tres dimensiones independientes, la estabilidad, la plausibilidad y la fidelidad, que ningún método optimiza a la vez y que la práctica debe ponderar según el propósito de cada auditoría. Este desplazamiento, desde la búsqueda del mejor explicador hacia la comprensión del compromiso entre dimensiones, es quizá el aporte más duradero del trabajo, ya que ofrece un marco de decisión en lugar de una recomendación cerrada.

El estudio deja también una lección metodológica que conviene subrayar. La disposición a corregir el rumbo cuando la evidencia contradice una hipótesis de partida, y a retractar hallazgos apoyados en artefactos de cómputo, no debilitó las conclusiones sino que las volvió más sólidas. En un campo joven como la explicabilidad de las redes neuronales sobre grafos, donde los resultados son sensibles al protocolo y las limitaciones de los datos se confunden con facilidad con las de los métodos, esta actitud de honestidad metodológica es, más que una virtud, una necesidad. Los hallazgos que aquí se reportan aspiran a ser útiles precisamente por haber sido sometidos a esa exigencia, y por distinguir con cuidado entre lo que la evidencia respalda y lo que apenas sugiere.
